% OVERLEAF COPY-PASTE VERSION
%
% Required packages in the main document preamble:
%   \usepackage{graphicx}
%   \usepackage{booktabs}
%
% Upload these three files from report/generated/ to the Overleaf project:
%   main_performance_vs_time.pdf
%   full_data_training_curves.pdf
%   robustness_curves.pdf
%
% Copy each marked block into the matching part of the group's CVPR document.
% Do not add a second \begin{document}, Abstract, Discussion, or Conclusion if
% the main document already contains one. All tables are inline, so this file
% does not depend on any external table include files.
%
% Quantitative claims are recorded in docs/member_e_claim_evidence.md.

% ======================== PASTE INTO ABSTRACT ========================
\begin{abstract}
Fine-grained species recognition is difficult when classes differ in small
morphological details and labelled examples are scarce. We evaluate a fixed
500-class subset of iNaturalist 2021 under a shared data split and evaluation
protocol. With 40 training images per class, we compare HOG and SIFT
pipelines, a ResNet18 trained from random initialization, and frozen or fully
fine-tuned ImageNet representations. We also study augmentation strength,
regularization, added self-attention, and visual explanations. Full
fine-tuning reaches 66.24\% Top-1 accuracy and 66.20\% Macro-F1, exceeding
the best scratch model by 28.30 and 28.69 percentage points, respectively.
Basic augmentation works better than either no augmentation or a stronger
policy, while randomly initialized self-attention does not improve
generalization. In a separate extension, we increase the cleaned training set
from 20,000 to 129,378 images while keeping the validation and test sets
fixed. Under the same 15-epoch budget, fine-tuning reaches 81.58\% Top-1,
compared with 64.20\% for scratch training. Pretrained representations
therefore remain useful as the training set grows, but augmentation strength
and model capacity still need to match the available supervision.
\end{abstract}

% ======================== PASTE INTO METHODS ========================
\subsection{Evaluation and Analysis Protocol}
\label{sec:evaluation_protocol}

Every method is evaluated on the same held-out set of 5,000 images, with ten
images from each of 500 classes. Model selection and hyperparameter choices
use only the validation split. The test split is reserved for final
evaluation, fixed-model robustness tests, and error analysis. We report Top-1
and Top-5 accuracy, macro-averaged precision, recall and F1, and balanced
accuracy. The macro metrics weight every class equally and complement
sample-level accuracy.

Rather than copy metrics from training logs, we recompute them from each
method's saved score matrix. Before generating the tables, the evaluation
program checks that every run contains exactly 5,000 unique image identifiers,
uses the shared 500-class score mapping, and assigns the same ground-truth
label to each image.

Training and inference times provide context, not a controlled hardware
benchmark. The limited-data methods were run on different CPU and GPU systems.
The two full-data jobs used the same NVIDIA L40S GPU model and software
environment, but they ran on different physical nodes, so their elapsed times
are only approximately comparable.

For the robustness test, four degradations are applied to decoded RGB images
before model-specific preprocessing: additive Gaussian noise, Gaussian blur,
brightness reduction, and JPEG compression. Each degradation has five
severity levels, and model parameters remain fixed throughout. We report the
score at each level and its change from an undegraded severity-0 run in the
same environment. Grad-CAM is used separately for qualitative failure
analysis. Its heat maps show where predictions are sensitive, but they do not
establish causal explanations.

% ================= PASTE INTO EXPERIMENTAL RESULTS =================
\subsection{Main Comparison}
\label{sec:main_comparison}

Table~\ref{tab:main_results} compares the method families in the limited-data
setting. The handcrafted baselines struggle with the 500-class task:
SIFT--BoVW with an SVM reaches 2.94\% Top-1, and HOG with a random forest
reaches 1.38\%. Training ResNet18 from scratch raises Top-1 to 37.94\%, showing
that learned visual features are considerably more effective even with
limited labels.

ImageNet pretraining improves the result further. Freezing the backbone and
training only the classifier gives 53.88\% Top-1, 15.94 points above the best
scratch run. Fine-tuning the full network reaches 66.24\% Top-1 and 66.20\%
Macro-F1. This adds 12.36 points over the frozen strategy and 28.30 points
over scratch training. Accuracy, balanced accuracy, and Macro-F1 are close
because the test split contains exactly ten images per class.

\begin{table}[t]
  \centering
  \caption{Limited-data results on the shared 5,000-image test set. Values
  other than time are percentages. Runtimes were measured on different
  hardware and should not be compared as a controlled speed benchmark.}
  \label{tab:main_results}
  \resizebox{\linewidth}{!}{%
    \begin{tabular}{lrrrrrr}
      \toprule
      Method & Top-1 & Top-5 & Macro-F1 & Bal. Acc. & Train (s) & Infer (s) \\
      \midrule
      SIFT + BoVW + SVM & 2.94 & 8.02 & 2.26 & 2.94 & 117.5 & 92.7 \\
      HOG + Random Forest & 1.38 & 5.26 & 0.89 & 1.38 & 686.7 & 37.3 \\
      Scratch (basic aug., 100-epoch SGD) & 37.94 & 61.24 & 37.51 & 37.94 & 4330.7 & 22.6 \\
      Pretrained ResNet18 (frozen) & \underline{53.88} & \underline{76.68} & \underline{53.55} & \underline{53.88} & 557.8 & 13.9 \\
      Pretrained ResNet18 (fine-tuned) & \textbf{66.24} & \textbf{85.82} & \textbf{66.20} & \textbf{66.24} & 725.7 & 12.8 \\
      \bottomrule
    \end{tabular}%
  }
\end{table}

Figure~\ref{fig:performance_time} places accuracy beside the recorded training
and inference times. It should not be read as a controlled speed comparison:
the CPU traditional pipelines and GPU ResNet18 runs used different systems.
The useful comparison is the large difference in representation quality, not
a platform-specific runtime ranking.

\begin{figure}[t]
  \centering
  \includegraphics[width=\linewidth]{main_performance_vs_time.pdf}
  \caption{Top-1 accuracy and recorded training and inference time for the
  main methods. Hardware labels are included because runtime was not
  controlled across the limited-data experiments.}
  \label{fig:performance_time}
\end{figure}

\subsection{Full-Data Scaling Extension}
\label{sec:full_data_results}

The full-data experiment tests whether ImageNet initialization still helps
under a matched training configuration. Both ResNet18 models use the cleaned
129,378-image training set, fixed validation and test splits, basic
augmentation, AdamW with cosine scheduling, a batch size of 64, and a
15-epoch budget. The scratch model reaches 64.20\% Top-1 and 64.02\%
Macro-F1. ImageNet initialization with full fine-tuning raises these scores to
81.58\% and 81.53\%, gains of 17.38 and 17.52 points. Top-5 accuracy rises
from 84.76\% to 93.70\%.

\begin{table}[t]
  \centering
  \caption{Full-data results under a matched 15-epoch ResNet18 training
  budget. Values other than time are percentages.}
  \label{tab:full_data_results}
  \resizebox{\linewidth}{!}{%
    \begin{tabular}{lrrrrrr}
      \toprule
      Method & Top-1 & Top-5 & Macro-F1 & Bal. Acc. & Train (s) & Infer (s) \\
      \midrule
      Full-data Scratch ResNet18 & \underline{64.20} & \underline{84.76} & \underline{64.02} & \underline{64.20} & 1295.0 & 13.7 \\
      Full-data Fine-tuned ResNet18 & \textbf{81.58} & \textbf{93.70} & \textbf{81.53} & \textbf{81.58} & 1351.7 & 5.3 \\
      \bottomrule
    \end{tabular}%
  }
\end{table}

The fine-tuned job takes 1,351.7 seconds, compared with 1,295.0 seconds for
scratch training, a difference of 56.7 seconds (4.38\%). Both nodes had an
L40S GPU, but they were different machines, so the timing difference is not a
strict measure of fine-tuning overhead. The best validation accuracies occur
at epochs 15 and 13 for scratch and fine-tuning, respectively.
Figure~\ref{fig:full_data_curves} shows stronger validation performance for
the fine-tuned model throughout training.

\begin{figure}[t]
  \centering
  \includegraphics[width=\linewidth]{full_data_training_curves.pdf}
  \caption{Training and validation curves for the matched full-data scratch
  and fine-tuned ResNet18 runs.}
  \label{fig:full_data_curves}
\end{figure}

The increase from the best limited-data scratch result (37.94\%) to the
full-data scratch result (64.20\%) is descriptive, not a pure data-scale
ablation. The limited-data run uses 100 epochs and SGD, whereas the full-data
run uses 15 epochs and AdamW. The matched full-data pair supports a narrower
claim: ImageNet initialization retains a substantial advantage at the larger
training scale.

\subsection{Robustness under Image Degradation}
\label{sec:robustness_results}

Figure~\ref{fig:robustness} summarizes all 20 corrupted test settings. The
fine-tuned model has the highest absolute Top-1 accuracy in every setting,
although the margin depends strongly on the corruption. Blur causes the
largest loss. At severity 5, fine-tuned Top-1 falls from 66.24\% to 9.02\%,
while scratch falls from 37.94\% to 3.90\%. Brightness reduction is less
damaging to the fine-tuned model: at factor 0.5 it retains 59.92\% Top-1,
compared with 16.70\% for scratch. Under the strongest Gaussian noise and
JPEG settings, the two models reach 44.24\% versus 27.70\%, and 50.12\%
versus 27.98\%, respectively.

\begin{figure}[t]
  \centering
  \includegraphics[width=\linewidth]{robustness_curves.pdf}
  \caption{Top-1 accuracy under four test-time degradations. Severity 0 is
  undegraded inference in the same local environment as severities 1--5.}
  \label{fig:robustness}
\end{figure}

Relative retention gives a slightly different view. At maximum blur,
brightness reduction, noise, and JPEG compression, fine-tuning retains
13.6\%, 90.5\%, 66.8\%, and 75.7\% of its clean accuracy. Scratch retains
10.3\%, 44.0\%, 73.0\%, and 73.7\%. After normalization by its lower clean
baseline, scratch is slightly better only under maximum Gaussian noise. The
small absolute drops for HOG and SIFT should not be mistaken for strong
robustness because both start close to the 0.2\% chance level. Noise is most
damaging to HOG (1.32\% to 0.60\%), while blur is most damaging to SIFT
(3.18\% to 1.26\%).

% ======================== PASTE INTO DISCUSSION ========================
\section{Discussion}
\label{sec:discussion}

The limited-data results favor learned and pretrained features. Handcrafted
descriptors do not capture enough of the color, texture, shape, and context
needed to separate 500 species. A scratch ResNet18 learns a stronger
representation, but 40 images per class are not enough for it to approach the
pretrained models. The frozen result shows that ImageNet features transfer to
this task, and the additional 12.36-point gain from fine-tuning shows the
benefit of adapting those features to species-specific cues.

More augmentation or capacity did not always help. Basic augmentation raises
scratch Top-1 from 22.62\% to 31.76\%, while the stronger policy reaches only
25.88\%. One likely explanation is that aggressive transformations disturb
color, texture, or local morphology that separates related species.
Regularization reduces the reported train--validation gap with only a small
Top-1 cost, but adding self-attention lowers Top-1 from 66.24\% to 62.08\%.
The randomly initialized attention parameters may add more capacity than the
limited training set can support. Because we tested only one attention
configuration, this result does not rule out other attention architectures.

Absolute robustness and relative retention tell slightly different stories.
Fine-tuning remains the most accurate method under every tested corruption
and is particularly stable to brightness reduction. Scratch, however,
preserves a slightly larger fraction of its clean accuracy under strong
Gaussian noise. Blur removes spatial detail used by both deep models and
causes their steepest decline. The flatter traditional curves do not indicate
strong robustness; HOG and SIFT already begin near the 0.2\% chance level.

The full-data extension narrows the gap between scratch training and
fine-tuning but does not remove it. Under the matched 15-epoch configuration,
fine-tuning still leads by 17.38 Top-1 points. This is consistent with scratch
training benefiting from more supervision. However, the limited-data and
full-data scratch configurations use different optimizers and epoch budgets,
so the change in their gap is not a controlled scaling law.

The study has several limits. It uses one balanced 500-class subset of
iNaturalist 2021, one main backbone, and generally one training seed per
configuration. We therefore cannot measure run-to-run variance or assume that
the ranking will hold for other datasets and architectures. Limited-data
runtimes also come from different hardware. The full-data comparison covers
only two initialization strategies under one 15-epoch budget, with the jobs
running on different L40S nodes. Finally, Grad-CAM is a qualitative
localization tool rather than evidence of causal reasoning. Future work
should repeat the main runs across several seeds and test modern fine-grained
backbones, taxonomic supervision, and additional matched full-data variants.

% ======================== PASTE INTO CONCLUSION ========================
\section{Conclusion}
\label{sec:conclusion}

Under the shared 500-class protocol, pretrained ResNet18 representations
outperform handcrafted features and scratch training when only 40 labelled
images are available per class. Full fine-tuning reaches 66.24\% Top-1 in
this setting. The ablations also show that stronger perturbations and extra
randomly initialized modules do not necessarily improve generalization. In
the separate 129,378-image extension, fine-tuning reaches 81.58\% Top-1,
compared with 64.20\% for scratch training under the matched configuration.
Pretraining therefore remains useful at the larger scale.

Fine-tuning is also the most accurate method under every tested degradation,
although blur exposes a shared dependence on fine spatial detail. Taken
together, the results support choosing representation, augmentation strength,
and model capacity according to the amount of available supervision, while
robustness tests and visual failure analysis show where the models still
break down.
